\documentclass[reprint,amsmath,amssymb,aps,twoside]{revtex4-2} % for editor use only

% For initial submission use these lines
%\documentclass[amsmath,amssymb,aps,endfloats,linenumbers]{revtex4-2}



% DO NOT CHANGE THINGS BELOW HERE
\usepackage{graphicx}
\usepackage{amsmath,amssymb,amsfonts}
\usepackage{dcolumn}
\usepackage{bm}
\usepackage{siunitx}
%\usepackage{tikz,pgfplots}
\sisetup{separate-uncertainty=true, multi-part-units=single}
\usepackage[colorlinks,allcolors=blue]{hyperref}
\usepackage{cleveref}
\crefname{equation}{}{}
\crefname{figure}{Fig.}{Figs.}
\crefname{table}{Table}{Tables}
\usepackage{svg}
\svgpath{{./figures}}
\makeatletter
\def\blfootnote{\xdef\@thefnmark{}\@footnotetext}
\makeatother





% set PDF metadata
% AUTHORS EDIT THIS PART
\hypersetup{%
pdftitle={Investigation of energy conservation using poppers}, % full title here
pdfauthor={Nimya Badmin, Kyle Koping, and Ella Pechersky}, % full author list here
}


\usepackage{fancyhdr}
\pagestyle{fancy}
\fancyhf{}
\fancyhead[RE,RO]{J S\&E \textbf{2}, 102--105 (2026)} % FOR EDITOR ONLY
\fancyhead[LO]{Badmin, Koping, \& Pechersky} % if many authors
\fancyhead[LE]{Investigation of energy conservation using poppers} 
\fancyfoot[C]{\thepage}
\fancypagestyle{mytitlepage}{
\fancyhf{}
\fancyhead[C]{Journal of Science \& Engineering \textbf{2}, 102--105 (2026)} % FOR EDITOR ONLY
\fancyfoot[C]{\thepage}
%\fancyfoot[L]{\scriptsize\copyright\ 2026 Journal of Science \& Engineering, \href{https://creativecommons.org/licenses/by-nc/4.0/}{CC-BY-NC 4.0}}
\fancyfoot[L]{\scriptsize\copyright\ 2026 J S\&E \href{https://creativecommons.org/licenses/by-nc/4.0/}{CC-BY-NC 4.0}}
}


\begin{document}
\setcounter{page}{102} % FOR EDITOR USE ONLY
%\linenumbers % FOR EDITOR USE ONLY

% AUTHORS EDIT THIS PART
\title{Investigation of energy conservation using poppers}
\author{Nimya Badmin}
\author{Kyle Koping}
\email[Contact author: ]{427kkoping@frhsd.com} % use FRHSD address
\author{Ella Pechersky}
% For S&E authors
\affiliation{Science \& Engineering Magnet Program, \href{https://manalapan.frhsd.com/}{Manalapan High School}, Englishtown, NJ 07726 USA}
%\date{\today}
\received{30 January 2026}
\accepted[accepted~]{11 April 2026}
\published[published~]{17 May 2026}


% AUTHORS PUT YOUR MANUSCRIPT IN BELOW HERE
\begin{abstract}
Rudolf Clausius proposed that energy cannot be created or destroyed, only transformed. This proposal is the basis of the first law of thermodynamics, the conservation of energy. This law indicates that for an object undergoing vertical motion, the initial kinetic energy is equal to the gravitational potential energy at the highest point in an ideal system. Therefore, we tested the validity of his claim. We measured a popper toy’s velocity immediately after launching and at its maximum height using video digitization. We then calculated its mechanical energy to test whether energy was conserved throughout the launch. We carried out a one-sample $t$-test using the differences in mechanical energy to determine whether the energies at both flight points were equal. The analysis provided no statistically significant evidence that the mechanical energy of the popper was not conserved, which is consistent with Clausius’ theory of conservation of energy.

\vspace{1em} % Add space before DOI
\noindent DOI: \href{https://doi.org/10.64804/mjhvfp73}{10.64808/mjhvfp73}
\end{abstract}

% For revtex format this will not print on the article page but will be used in setting 
% the journal metadata so include keywords. Lowercase unless they are acronyms or proper nouns
% and pick keywords that will actually help people find your work properly. 
\keywords{conservation of energy, potential energy, kinetic energy, velocity, maximum height, peak height, FizziQ, digitization, mean, standard deviation, t-test, error}

\maketitle\thispagestyle{mytitlepage}
\blfootnote{\\\emph{Published by the Journal of Science \& Engineering under the terms of the \href{https://creativecommons.org/licenses/by-nc/4.0/}{Creative Commons Attribution-NonCommercial 4.0 International} license. Further distribution of this work must maintain attribution to the author(s) and the published article's title, journal citation, and DOI.}}

%The current column width is \the\columnwidth






\section{Introduction}
The law of conservation of energy states that energy cannot be created or destroyed, but can only be transformed from one form to another, such as from potential to kinetic energy \cite{clausius:1851:on,tipler,barrons,openstax}. In an ideal system with no non-conservative forces, the mechanical energy of the system remains constant.

Potential energy is the energy stored in an object relative to its position \cite{tipler,barrons,openstax}:
\begin{equation}
\text{PE} = mgh.
\end{equation}

Kinetic energy is the energy of motion \cite{tipler,barrons,openstax}:
\begin{equation}
\text{KE} = \frac{1}{2} m v^2.
\end{equation}

Total mechanical energy is the sum of a system’s potential and kinetic energy \cite{tipler,barrons,openstax}:
\begin{equation}
\text{ME} = \text{PE} + \text{KE}.
\end{equation}

The purpose of this experiment was to test the principle of energy conservation by calculating and comparing the popper’s total mechanical energy at launch and at its peak height.

The null hypothesis ($H_0$) states that the mean difference in the popper’s mechanical energy at launch and at peak height is zero, meaning that the total mechanical energy remains constant and energy is conserved:
\begin{equation}
H_0: E(\text{ME}_f - \text{ME}_0) = 0.
\end{equation}
The alternative hypothesis ($H_a$) states that the mean difference in the popper’s kinetic energy at launch and at peak height is positive, meaning that mechanical energy is greater at launch than at peak height, indicating that some mechanical energy is lost during flight, and therefore, energy is not conserved:
\begin{equation}
H_a: E(\text{ME}_f - \text{ME}_0) > 0.
\end{equation}
    
    
    
    
    



\section{Methods and materials}
\subsection{Lab materials and setup}
This experiment included the use of a popper (Jumping Emoticon Popper Spring Launchers; Liberty Imports; Allentown, PA), a scissor jack, a digital scale (TOP2KG; SmartWeight; Jiangsu, China), and a meter stick, as shown in in \cref{fig:212223}.
\begin{figure}
\begin{center}
\includegraphics[width=\columnwidth]{figures/setup.png}
\end{center}
\caption{(A) Uncompressed popper. (B) Compressed popper. (C) Setup for measuring mechanical work using digital scale and scissor jack.}
\label{fig:212223} 
\end{figure}
To examine energy conservation, we launched a popper with a height of approximately 3.5-inch (\qty{89}{\milli\meter}; see \cref{fig:212223}A) from its maximum possible compression of approximately 2-inch (\qty{51}{\milli\meter}; see \cref{fig:212223}B).  For digitization, a smartphone camera (iPhone 13, recording at 60fps at 1.0x zoom in HD; Apple Inc; Cupertino, CA) and video analysis software (FizziQ, version 5.0.4) were used \cite{chazot:2025:using,badmin:2026:investigation}. In FizziQ, digitization was performed at regular intervals (approximately every 6 frames), resulting in a total of 4-6 data points per trial, similar to \cite{badmin:2026:investigation}. A meter stick was set up vertically for scale, to measure the maximum height attained by the popper during launch. The popper was set approximately 2-inch from the yardstick so that the yardstick would not interfere with the experiment.

The mass of the popper was \qty{6.0}{\gram}, and the force required to fully compress the popper was \qty{778}{\gram}, which is \qty{7.62}{\newton}. The force was found by compressing the popper with a scissor jack on a digital scale  (see \cref{fig:212223}C); this method was also used to determine work done in compressing the popper, as described below. 
    
\subsection{Method of testing and obtaining data}
For the data collection process, one person recorded the launches while another person launched the popper. The recorder was positioned so that the popper’s peak height was captured on video. The recorder would begin recording and count down from five seconds before launch. The launcher would first compress the popper by pushing down on it, and then release their hand after the countdown to launch the popper into the air. This process was repeated five times to account for confounding variables, including delayed countdowns, human reaction time, and poor launching technique.

These videos were then digitized using FizziQ to find the launch velocity at the earliest possible point \cite{chazot:2025:using,badmin:2026:investigation}. The scale of the video was set using the meter stick as a scale reference. Then, a single point on the popper was tracked throughout its motion, with data points recorded at regular time intervals (approximately every 6 frames). These points were then used by FizziQ to generate the position and velocity data shown in the raw results.

The results from the five launches were averaged before carrying out a one-sample $t$-test. Calculations for the statistical test were performed using a TI-84 graphing calculator. Additional statistical analyses including ANOVA were performed in R \cite{starnes:2015:practice,R,ggplot2,dplyr}. 

To determine work done in compressing the popper, the popper was placed on the digital scale and the scissor jack was used to apply a controlled force during compression. A meter stick was placed beside the scale to measure the displacement of the popper as it was compressed. The data was recorded at multiple intervals (approximately every \qty{0.005}{\meter} of compression), with  three replicates for each compression value.
    
Data and analysis code are available at \url{https://github.com/devangel77b/427kkoping-lab3}.    
    
    
    
    
    
    
    
\section{Results}
The mass of the popper was measured to be \qty{6.0}{\gram} (\qty{0.0060}{\kilo\gram}). The force required to fully compress the popper was measured as \qty{778}{\gram} force, which is equivalent to \qty{7.62}{\newton}, using $F = mg$ where $g=\qty{9.8}{\meter\per\second\squared}$ \cite{tipler,barrons,openstax}.

% You should use something more appropriate like R
\begin{table}
\caption{Popper initial velocity, height, KE and GPE}
\label{tab:3132}
\begin{center}
\begin{ruledtabular}
\begin{tabular}{SSSSS}
{trial} & {$v_0$, \unit{\meter\per\second}} & {$h$, \unit{\meter}} & {KE, \unit{\joule}} & {GPE, \unit{\joule}} \\
\colrule \\
1 & 5.29 & 1.43 & 0.0840 & 0.0842 \\
2 & 5.27 & 1.45 & 0.0833 & 0.0853 \\
3 & 5.06 & 1.33 & 0.0768 & 0.0783 \\
4 & 5.50 & 1.67 & 0.0908 & 0.0983 \\
5 & 5.68 & 1.54 & 0.0968 & 0.0906 \\
\end{tabular}
\end{ruledtabular}
\end{center}
\end{table}

\begin{figure}
\begin{center}
\includesvg[width=\columnwidth]{fig33}
\end{center}
\caption{Bar graph of kinetic and potential energies; data from \cref{tab:3132}. KE and PE are not significantly different (ANOVA; $p=0.836$).}
\label{fig:33}
\end{figure}

\begin{table}
\caption{Force-compression raw data for \cref{fig:35}, to determine work to compress a popper}
\label{tab:34}
\begin{center}
\begin{ruledtabular}
\begin{tabular}{S[table-format=1.3]*3S[table-format=1.2]}
{$x$, \unit{\meter}} & \multicolumn{3}{c}{$F$, \unit{\newton}} \\
\colrule\\
0.000 & 0.00 & 0.00 & 0.00 \\
0.008 & 1.33 & 1.36 & 0.91 \\
0.013 & 2.39 & 2.00 & 1.98 \\
0.018 & 2.88 & 2.87 & 3.26 \\
0.023 & 3.95 & 3.81 & 3.79 \\
0.028 & 5.25 & 5.28 & 5.16 \\
0.033 & 6.56 & 6.67 & 6.77 \\
0.038 & 7.18 & 7.18 & 7.27 \\
\end{tabular}
\end{ruledtabular}
\end{center}
\end{table}

\begin{figure}
\begin{center}
\includesvg[width=\columnwidth]{fig35}
\end{center}
\caption{Force (\unit{\newton}) versus compression (\unit{\meter}); data from \cref{tab:34}. The work done to compress the popper (shaded area) is $W=\qty{0.130\pm0.001}{\joule}$.}
\label{fig:35}
\end{figure}



\section{Discussion}
\subsection{Energy is conserved during flight phase}
Our null hypothesis ($H_0$) was that the mean difference in the popper’s mechanical energy at launch and at peak height is zero, meaning that the total mechanical energy remains constant. Our alternative hypothesis ($H_a$) was that the mean difference in the popper’s mechanical energy at launch and at peak height is positive, indicating that energy is lost during flight. Our results from both a one-sample $t$-test as well as from analysis of variance (ANOVA) fail to reject the null hypothesis and do not support the alternative hypothesis.  The values of the test statistic and $p$-value obtained were 0.239 and 0.411, respectively, as can be seen in \cref{tab:3132,fig:33}. Since the $p$-value is greater than the tested significance level of 0.05, there is no statistically significant difference between the mechanical energies at launch and at peak height. Therefore, there is no significant evidence to conclude that energy is lost during the popper’s motion. These results are most consistent with our null hypothesis. 

Although individual trials show small differences between mechanical energies at launch and at peak height, these variations are likely due to experimental error. Overall, the statistical analysis supports the idea of energy conservation. However, this result is limited by the small sample size (five replicates; one popper) and potential sources of error.

%Overall, the $t$-test fails to reject our null hypothesis ($H_0$) that energy is conserved. 
%The graph of kinetic and potential energies (see \cref{tab:34}) shows a similar trend, with the exception of one obvious outlier (trial two), where the popper appears to have lost 18.01 mJ of energy. While the t-test and graph support energy conservation, individual trial results show variability rather than a consistent loss of energy, most likely due to experimental errors.

\subsection{Energy lost during launch is significant}
Kinetic energy at launch (\cref{fig:33}) is significantly less than the work done to compress the popper (\cref{fig:35}. With the data we obtained, we are unable to explain why the launch process loses approximately \qtyrange{30}{40}{\percent} of the elastically potential energy stored when the popper is compressed; this should be examined further in future work.

\subsection{Sources of error}
Many possible sources of error could have occurred during this experiment. For example, air resistance could cause some of the popper’s mechanical energy to be lost as thermal energy and sound, resulting in slightly lower mechanical energy at peak height, although \cref{fig:33} shows no significant differences between KE and GPE.

Another possible source of error could be inaccuracies in digitizing the popper’s motion. Camera angle distortion, limited frame rate, and human error in placing tracking points could all affect the measured position and velocity value. Additionally, variations in the popper release technique may have affected the results. Differences in how the popper was compressed and released could lead to inconsistent initial energy between trials.

%\subsection{Conclusion}



    
    
\section{Acknowledgements}
We thank the Science \& Engineering Program at Manalapan High School for support. NB helped with methods and materials and results.  KK completed the abstract, carried out the t-test, and wrote about its results. EP helped with the launch procedure, did the video motion analysis, worked on introduction, methods and materials, and discussion, and handled initial and final revisions.


% References – Do not spend too much time trying to perfectly format per MLA or whatever because journals often have their own preferred formats; as long as you have author, title, journal, year, volume, page number it will be findable. Peer reviewed journals or published work only; Food Girl’s Blog is not a peer reviewed journal; Wikipedia is not a peer reviewed journal; Webster’s Dictionary is not necessary in a scientific journal article. If I find fake references your paper will be rejected. 






\bibliography{lab.bib}

%Chazot, Christophe. FizziQ. 5.0.4, Trapèze.digital, 2025, Apple App Store.

%Clausius, R. “On The Moving Force of Heat, and the Laws Regarding the Nature of Heat Itself Which Are Deducible Therefrom.” THE LONDON, EDINBURGH and DUBLIN PHILOSOPHICAL MAGAZINE AND JOURNAL OF SCIENCE, June 1851, sites.pitt.edu/~jdnorton/teaching/2559_Therm_Stat_Mech/docs/Clausius%20Moving%20Force%20heat% 201851.pdf.

%P. A. Tipler and G. Mosca, Physics for Scientists and Engineers, 5th ed. (W H Freeman and Company, New York, 2004).

\end{document}
